\programchapter{Code de conduite [FR]}

Les organisateurs souhaitent que cette rencontre soit productive et agr\'eable
pour tout le monde, peu importe le genre, l'orientation sexuelle, le handicap,
l'apparence physique, la corpulence, l'origine, la nationalit\'e, la religion,
le stade de carri\`ere ou l'affiliation institutionnelle. Le harc\`element, la
discrimination, l'intimidation, l'exclusion et les repr\'esailles n'y sont pas
tol\'er\'es.

En pratique, on s'attend \`a un comportement professionnel, \`a des \`echanges
adapt\'es \`a un public professionnel vari\'e, et \`a un traitement respectueux des
autres participant\,e\,s. Les commentaires sexistes, racistes ou d\'enigrants,
les interruptions soutenues, les contacts physiques inappropri\'es, les avances
non sollicitées, les photos ou enregistrements sans consentement, et les propos
ou images \`a caract\`ere sexuel ne sont pas appropri\'es.

Toute personne \`a qui l'on demande de cesser un comportement inappropri\'e doit
s'y conformer imm\'ediatement. Les organisateurs se r\'eservent le droit de
prendre les mesures jug\'ees n\'ecessaires, y compris le retrait de la rencontre
sans remboursement.

Pour signaler une situation en toute confidentialit\'e, on peut s'adresser \`a
une personne de la pr\'esidence du LOC, \`a un membre du comit\'e EDI de la CASCA,
ou \`a un membre du conseil d'administration ou \`a un ombudsperson de la CASCA.

Le texte officiel faisant autorit\'e demeure celui de la CASCA; cette version
bref r\'esume l'approche utilis\'ee dans le programme de la r\'eunion annuelle 2025
et renvoie au document officiel ici :

\begin{itemize}
  \item \href{https://casca.ca/?page_id=22442}{Page CASCA: Code of Conduct}
  \item \href{https://casca.ca/wp-content/uploads/2025/12/The-CASCA-Code-of-Conduct.pdf}{Document officiel PDF}
\end{itemize}